\documentclass{article}
\usepackage{amsmath,amssymb,amsthm}
\usepackage{algorithm}
\usepackage{algorithmic}
\usepackage{bm}
\usepackage{enumitem}
\usepackage{hyperref}
\usepackage{geometry}
\geometry{margin=1in}
\newtheorem{theorem}{Theorem}
\newtheorem{lemma}{Lemma}
\newtheorem{assumption}{Assumption}
\newcommand{\EE}{\mathbb{E}}
\newcommand{\RR}{\mathbb{R}}
\newcommand{\norm}[1]{\left\lVert #1 \right\rVert}
\newcommand{\inner}[2]{\left\langle #1,#2 \right\rangle}
\newcommand{\prox}{\operatorname{prox}}
\begin{document}

\section*{Algorithm Name}
\textbf{Differentially Private Stochastic Gradient Descent with Gaussian Noise (DP‑SGD‑G)}  

The algorithm performs (full‑batch) gradient descent on a convex, $L$‑smooth objective while adding independent isotropic Gaussian noise to each gradient estimate in order to satisfy $(\varepsilon,\delta)$‑differential privacy.

\section*{Mathematical Formulation}
Let $f:\RR^{d}\to\RR$ be convex and $L$‑smooth.  
Given an initial point $x_{0}\in\RR^{d}$, a learning‑rate schedule $\{\eta_{t}\}_{t\ge 0}$, and a noise scale $\sigma>0$, the iterates are defined by  
\[
\boxed{%
x_{t+1}=x_{t}-\eta_{t}\bigl(\nabla f(x_{t})+\xi_{t}\bigr),\qquad 
\xi_{t}\sim\mathcal{N}\bigl(0,\sigma^{2}I_{d}\bigr)
}\tag{1}
\]
where the random vectors $\xi_{t}$ are i.i.d. and independent of the data.  

For privacy accounting we keep track of the per‑iteration privacy loss $\varepsilon_{t}$ of the Gaussian mechanism (see \cite{dwork2014algorithmic}); the total privacy budget after $T$ steps is obtained by standard composition (e.g., moments accountant).  

\section*{Pseudocode}
\begin{algorithm}[H]
\caption{DP‑SGD‑G}
\begin{algorithmic}[1]
\REQUIRE Convex $L$‑smooth $f$, initial point $x_{0}$, learning‑rate schedule $\{\eta_{t}\}$, noise scale $\sigma$, total iterations $T$
\STATE Set $x\gets x_{0}$
\FOR{$t=0$ \TO $T-1$}
    \STATE Compute exact gradient $g_{t}\gets \nabla f(x)$
    \STATE Sample $\xi_{t}\sim\mathcal{N}(0,\sigma^{2}I_{d})$
    \STATE Noisy gradient $\tilde g_{t}\gets g_{t}+\xi_{t}$
    \STATE Update $x\gets x-\eta_{t}\tilde g_{t}$
    \STATE (Optional) Update privacy accountant with $(\sigma,\eta_{t})$
\ENDFOR
\RETURN $x$
\end{algorithmic}
\end{algorithm}

\section*{Dry Run Example}
Consider the one‑dimensional quadratic
\[
f(w)=(w-3)^{2},\qquad \nabla f(w)=2(w-3).
\]
Set $w_{0}=0$, constant step size $\eta=0.1$, noise variance $\sigma^{2}=0.01$, and run three iterations.

\begin{center}
\begin{tabular}{c|c|c|c}
$t$ & $\nabla f(w_{t})$ & $\xi_{t}\sim\mathcal{N}(0,0.01)$ (sample) & $w_{t+1}=w_{t}-\eta(\nabla f(w_{t})+\xi_{t})$\\ \hline
0 & $2(0-3)=-6$ & $0.05$ & $w_{1}=0-0.1(-6+0.05)=0.595$\\
1 & $2(0.595-3)=-4.81$ & $-0.08$ & $w_{2}=0.595-0.1(-4.81-0.08)=1.083$\\
2 & $2(1.083-3)=-3.834$ & $0.02$ & $w_{3}=1.083-0.1(-3.834+0.02)=1.461$\\
\end{tabular}
\end{center}
Objective values: $f(w_{0})=9$, $f(w_{1})\approx5.84$, $f(w_{2})\approx3.75$, $f(w_{3})\approx2.31$.  
The noisy updates still decrease the objective on average.

\newpage
\section*{Assumptions}
\begin{assumption}[Convexity and Smoothness]\label{as:convex}
$f$ is convex and $L$‑smooth, i.e.
\[
\forall x,y\in\RR^{d}:\qquad 
f(y)\le f(x)+\inner{\nabla f(x)}{y-x}+\frac{L}{2}\norm{y-x}^{2}.
\tag{A1}
\]
\end{assumption}
\begin{assumption}[Bounded Gradient Norm]\label{as:bounded}
There exists $G>0$ such that $\norm{\nabla f(x)}\le G$ for all $x$ in the domain visited by the algorithm.
\tag{A2}
\end{assumption}
\begin{assumption}[Noise Distribution]\label{as:noise}
The added noise satisfies $\xi_{t}\stackrel{\text{i.i.d.}}{\sim}\mathcal{N}(0,\sigma^{2}I_{d})$ and is independent of $x_{t}$.
\tag{A3}
\end{assumption}
\begin{assumption}[Learning‑Rate]\label{as:lr}
The step size is constant $\eta_{t}\equiv\eta\in(0,\tfrac{1}{L}]$.
\tag{A4}
\end{assumption}

\section*{Main Convergence Theorem}
\begin{theorem}\label{thm:conv}
Under Assumptions \ref{as:convex}–\ref{as:lr}, let $x^{\ast}\in\arg\min_{x}f(x)$ and define the averaged iterate $\bar x_{T}\!=\!\tfrac{1}{T}\sum_{t=0}^{T-1}x_{t}$. Then
\[
\EE\bigl[f(\bar x_{T})-f(x^{\ast})\bigr]
\le
\frac{\norm{x_{0}-x^{\ast}}^{2}}{2\eta T}
+\frac{\eta}{2}\bigl(G^{2}+d\sigma^{2}\bigr).
\tag{2}
\]
Consequently, choosing $\eta=\Theta\!\bigl(\tfrac{1}{\sqrt{T}}\bigr)$ yields
\[
\EE\bigl[f(\bar x_{T})-f(x^{\ast})\bigr]=\mathcal{O}\!\Bigl(\frac{1}{\sqrt{T}}\Bigr)
+ \mathcal{O}\!\Bigl(\frac{d\sigma^{2}}{\sqrt{T}}\Bigr).
\tag{3}
\]
\end{theorem}

\section*{Theoretical Properties}
\begin{itemize}[leftmargin=*]
\item \textbf{Stability.} The recursion (1) is a contraction in expectation as long as $\eta\le \tfrac{1}{L}$; the additional Gaussian term adds a bounded variance term $d\sigma^{2}$.
\item \textbf{Hyper‑parameter Sensitivity.} The bound (2) is linear in $\eta$ for the variance term and inversely proportional to $\eta$ for the initial‑error term, giving the classic trade‑off that leads to the optimal $\eta\asymp 1/\sqrt{T}$.
\item \textbf{Comparison to Vanilla SGD.} Setting $\sigma=0$ recovers the standard SGD bound $\frac{\norm{x_{0}-x^{\ast}}^{2}}{2\eta T}+\frac{\eta G^{2}}{2}$.
\item \textbf{Privacy‑Utility Trade‑off.} Larger $\sigma$ improves privacy (smaller $\varepsilon$) but degrades the term $d\sigma^{2}$ in (2), reflecting the well‑known utility loss due to noise.
\end{itemize}

\section*{Discussion}
The theorem shows that DP‑SGD‑G attains the same $O(1/\sqrt{T})$ convergence rate as non‑private SGD up to an additive term proportional to $d\sigma^{2}$. This term is unavoidable for Gaussian mechanisms \cite{dwork2014algorithmic}. By calibrating $\sigma$ according to a target $(\varepsilon,\delta)$ budget (e.g., via the moments accountant), one can obtain explicit utility guarantees for any fixed privacy level.

\newpage
\section*{Preliminaries (Definitions and Lemmas)}
\begin{lemma}[Three‑Point Identity for $L$‑smooth functions]\label{lem:three}
For any $x,y\in\RR^{d}$,
\[
f(y)\le f(x)+\inner{\nabla f(x)}{y-x}+\frac{L}{2}\norm{y-x}^{2}.
\tag{L1}
\]
\end{lemma}
\begin{lemma}[Unbiasedness of the noise]\label{lem:noise}
For every $t$, $\EE[\xi_{t}]=0$ and $\EE[\norm{\xi_{t}}^{2}]=d\sigma^{2}$.
\tag{L2}
\end{lemma}
\begin{proof}
Since $\xi_{t}\sim\mathcal{N}(0,\sigma^{2}I_{d})$, the mean vector is zero and each coordinate has variance $\sigma^{2}$, giving $\EE[\norm{\xi_{t}}^{2}]=\sum_{i=1}^{d}\EE[\xi_{t,i}^{2}]=d\sigma^{2}$.
\end{proof}

\section*{Main Proof (Step‑by‑Step Derivation)}
We prove Theorem \ref{thm:conv}. All expectations are taken over the randomness of $\{\xi_{t}\}_{t\ge0}$.

\begin{proof}
\textbf{Step 1 (Expand squared distance).}  
For any $t\ge0$,
\[
\begin{aligned}
\norm{x_{t+1}-x^{\ast}}^{2}
&=\norm{x_{t}-\eta(\nabla f(x_{t})+\xi_{t})-x^{\ast}}^{2}\\
&=\norm{x_{t}-x^{\ast}}^{2}
-2\eta\inner{\nabla f(x_{t})+\xi_{t}}{x_{t}-x^{\ast}}
+\eta^{2}\norm{\nabla f(x_{t})+\xi_{t}}^{2}. \tag{4}
\end{aligned}
\]
\textbf{Step 2 (Take expectation, use linearity).}  
Condition on $x_{t}$ and apply Lemma \ref{lem:noise}:
\[
\EE\!\bigl[\norm{x_{t+1}-x^{\ast}}^{2}\mid x_{t}\bigr]
=\norm{x_{t}-x^{\ast}}^{2}
-2\eta\inner{\nabla f(x_{t})}{x_{t}-x^{\ast}}
+\eta^{2}\EE\!\bigl[\norm{\nabla f(x_{t})+\xi_{t}}^{2}\mid x_{t}\bigr]. \tag{5}
\]
(The cross term $\inner{\xi_{t}}{x_{t}-x^{\ast}}$ vanishes because $\EE[\xi_{t}]=0$.)

\textbf{Step 3 (Expand the norm inside the expectation).}
\[
\begin{aligned}
\EE\!\bigl[\norm{\nabla f(x_{t})+\xi_{t}}^{2}\mid x_{t}\bigr]
&=\norm{\nabla f(x_{t})}^{2}
+2\inner{\nabla f(x_{t})}{\EE[\xi_{t}]}
+\EE\!\bigl[\norm{\xi_{t}}^{2}\bigr] \\
&=\norm{\nabla f(x_{t})}^{2}+d\sigma^{2}. \tag{6}
\end{aligned}
\]

\textbf{Step 4 (Insert (6) into (5)).}
\[
\EE\!\bigl[\norm{x_{t+1}-x^{\ast}}^{2}\mid x_{t}\bigr]
=\norm{x_{t}-x^{\ast}}^{2}
-2\eta\inner{\nabla f(x_{t})}{x_{t}-x^{\ast}}
+\eta^{2}\norm{\nabla f(x_{t})}^{2}
+\eta^{2}d\sigma^{2}. \tag{7}
\]

\textbf{Step 5 (Apply convexity).}  
By convexity of $f$,
\[
f(x_{t})-f(x^{\ast})\le \inner{\nabla f(x_{t})}{x_{t}-x^{\ast}}. \tag{8}
\]
Rearranging gives $\inner{\nabla f(x_{t})}{x_{t}-x^{\ast}}\ge f(x_{t})-f(x^{\ast})$.

\textbf{Step 6 (Upper‑bound the quadratic term).}  
Using $2ab\le a^{2}+b^{2}$ with $a=\eta\norm{\nabla f(x_{t})}$ and $b=\norm{x_{t}-x^{\ast}}$,
\[
-2\eta\inner{\nabla f(x_{t})}{x_{t}-x^{\ast}}
\le -2\eta\bigl(f(x_{t})-f(x^{\ast})\bigr). \tag{9}
\]

\textbf{Step 7 (Combine (7), (8), (9)).}
\[
\EE\!\bigl[\norm{x_{t+1}-x^{\ast}}^{2}\mid x_{t}\bigr]
\le
\norm{x_{t}-x^{\ast}}^{2}
-2\eta\bigl(f(x_{t})-f(x^{\ast})\bigr)
+\eta^{2}\norm{\nabla f(x_{t})}^{2}
+\eta^{2}d\sigma^{2}. \tag{10}
\]

\textbf{Step 8 (Bound the gradient norm).}  
Assumption \ref{as:bounded} gives $\norm{\nabla f(x_{t})}^{2}\le G^{2}$.

\textbf{Step 9 (Take full expectation).}
Denote $\Delta_{t}\coloneqq\EE\!\bigl[f(x_{t})-f(x^{\ast})\bigr]$. Taking expectation of (10) yields
\[
\EE\!\bigl[\norm{x_{t+1}-x^{\ast}}^{2}\bigr]
\le
\EE\!\bigl[\norm{x_{t}-x^{\ast}}^{2}\bigr]
-2\eta\Delta_{t}
+\eta^{2}\bigl(G^{2}+d\sigma^{2}\bigr). \tag{11}
\]

\textbf{Step 10 (Telescoping sum).}  
Sum (11) over $t=0,\dots,T-1$:
\[
\EE\!\bigl[\norm{x_{T}-x^{\ast}}^{2}\bigr]
\le
\norm{x_{0}-x^{\ast}}^{2}
-2\eta\sum_{t=0}^{T-1}\Delta_{t}
+T\eta^{2}\bigl(G^{2}+d\sigma^{2}\bigr). \tag{12}
\]

\textbf{Step 11 (Drop the non‑negative term).}  
Since the left‑hand side is non‑negative,
\[
2\eta\sum_{t=0}^{T-1}\Delta_{t}
\le
\norm{x_{0}-x^{\ast}}^{2}
+T\eta^{2}\bigl(G^{2}+d\sigma^{2}\bigr). \tag{13}
\]

\textbf{Step 12 (Average the gaps).}  
Define the averaged iterate $\bar x_{T}=\frac{1}{T}\sum_{t=0}^{T-1}x_{t}$. By convexity of $f$,
\[
f(\bar x_{T})-f(x^{\ast})
\le \frac{1}{T}\sum_{t=0}^{T-1}\bigl(f(x_{t})-f(x^{\ast})\bigr).
\]
Taking expectation and using (13):
\[
\EE\!\bigl[f(\bar x_{T})-f(x^{\ast})\bigr]
\le
\frac{1}{2\eta T}\norm{x_{0}-x^{\ast}}^{2}
+\frac{\eta}{2}\bigl(G^{2}+d\sigma^{2}\bigr). \tag{14}
\]

\textbf{Step 13 (Conclusion).}  
Equation (14) is exactly the bound (2) stated in Theorem \ref{thm:conv}. ∎
\end{proof}

\section*{Rate Derivation (With Constants)}
From (14) we have two competing terms:
\[
A(\eta)=\frac{\norm{x_{0}-x^{\ast}}^{2}}{2\eta T},
\qquad
B(\eta)=\frac{\eta}{2}\bigl(G^{2}+d\sigma^{2}\bigr).
\]
Minimising $A(\eta)+B(\eta)$ w.r.t. $\eta$ gives
\[
\eta^{\star}= \sqrt{\frac{\norm{x_{0}-x^{\ast}}^{2}}{T\bigl(G^{2}+d\sigma^{2}\bigr)}}
\le\frac{1}{L}\quad\text{(by assumption A4)}.
\]
Plugging $\eta^{\star}$ into (14) yields
\[
\EE\!\bigl[f(\bar x_{T})-f(x^{\ast})\bigr]
\le
\frac{\norm{x_{0}-x^{\ast}}}{\sqrt{T}}\sqrt{G^{2}+d\sigma^{2}}.
\tag{15}
\]
Thus the algorithm achieves an $O\!\bigl(\tfrac{1}{\sqrt{T}}\bigr)$ rate, with an explicit constant that grows as $\sqrt{d}\,\sigma$.

\section*{Special Cases}
\begin{enumerate}
\item \textbf{No privacy noise ($\sigma=0$).}  
Bound (2) reduces to the classical SGD bound:
\[
\EE[f(\bar x_{T})-f(x^{\ast})]\le
\frac{\norm{x_{0}-x^{\ast}}^{2}}{2\eta T}
+\frac{\eta G^{2}}{2}.
\]

\item \textbf{High‑dimensional regime.}  
When $d$ is large, the variance term $d\sigma^{2}$ dominates $G^{2}$, so the utility loss scales as $\mathcal{O}\!\bigl(\tfrac{d\sigma^{2}}{\sqrt{T}}\bigr)$. This matches known lower bounds for differentially private optimization \cite{bassily2014private}.

\item \textbf{Optimal step size under a fixed privacy budget.}  
If $\sigma$ is chosen to satisfy a given $(\varepsilon,\delta)$ budget (e.g., $\sigma=\frac{c\sqrt{T\log(1/\delta)}}{\varepsilon}$ for Gaussian mechanism), then (15) becomes
\[
\EE[f(\bar x_{T})-f(x^{\ast})]
= \mathcal{O}\!\Bigl(\frac{1}{\sqrt{T}}+\frac{d\,c^{2}\log(1/\delta)}{\varepsilon^{2}\sqrt{T}}\Bigr).
\]
\end{enumerate}

\newpage
\begin{thebibliography}{9}
\bibitem{dwork2014algorithmic}
C.~Dwork and A.~Roth, \emph{The algorithmic foundations of differential privacy}, Foundations and Trends in Theoretical Computer Science, 9(3–4):211–407, 2014.

\bibitem{bassily2014private}
R.~Bassily, A.~Smith, and A.~Thakurta, \emph{Private empirical risk minimization: Efficient algorithms and tight error bounds}, Proceedings of the 55th IEEE Symposium on Foundations of Computer Science (FOCS), 2014.
\end{thebibliography}

\end{document}